\section*{Chapitre \ref{ch:b3:genomics} --- Génomique et séquençage}

\begin{solution}{exo:b3:genomics:1}
Un \omterm{def:b3:genomics:sanger}{didésoxynucléotide} n'a pas d'hydroxyle 3$'$, donc aucune liaison
phosphodiester ne peut être formée avec le nucléotide suivant et la
chaîne s'arrête. Avec les seules formes didésoxy, toute chaîne
s'arrêterait à la première position~; avec les seules formes normales,
aucune ne s'arrêterait. Le mélange fait de la terminaison un événement
aléatoire à chaque position, si bien que les produits forment une
échelle complète, une longueur pour chaque base de la matrice.
\end{solution}

\begin{solution}{exo:b3:genomics:2}
$4\times 10^{8}\times \qty{300}{bp} = 1.2\times 10^{11}$ bases.
Homme~: $1.2\times 10^{11}/3.2\times 10^{9} \approx 38\times$.
Bactérie~: $1.2\times 10^{11}/5\times 10^{6} = \num{24000}\times$.
\end{solution}

\begin{solution}{exo:b3:genomics:3}
Un contig est une séquence continue assemblée à partir de \omterm{def:b3:genomics:ngs}{lectures}
chevauchantes~; un \omterm{def:b3:genomics:assembly}{échafaudage} est un jeu de contigs ordonnés et
orientés, séparés par des lacunes de taille estimée~; le \omterm{def:b3:genomics:assembly}{N50} est la
longueur de contig à laquelle les contigs triés atteignent la moitié
des bases assemblées. Ici les dix contigs de \qty{8}{Mb} contiennent à
eux seuls \qty{80}{Mb}, au-delà du seuil de \qty{50}{Mb} (le septième
atteint \qty{56}{Mb})~: \omterm{def:b3:genomics:assembly}{N50} $= \qty{8}{Mb}$.
\end{solution}

\begin{solution}{exo:b3:genomics:4}
La taille d'un génome ne suit pas la complexité de l'organisme~: un
oignon a cinq fois l'ADN d'un homme, un dipneuste quarante fois~; la
levure et le ver diffèrent d'un facteur dix en ADN pour des nombres de
gènes voisins. L'excédent est non codant~: éléments transposables et
leurs restes, introns, répétitions satellites.
\end{solution}

\begin{solution}{exo:b3:genomics:5}
$e^{-c} = 10^{-6}$ donne $c = 6\ln 10 = 13.8$. Pour $c = 8$~: $N =
cG/L = 8\times 5\times 10^{6}/150 \approx \num{267000}$ \omterm{def:b3:genomics:ngs}{lectures}, $\theta
= 30/150 = 0.2$, contigs $= N e^{-6.4} = \num{267000}\times 0.00166
\approx 440$.
\end{solution}

\begin{solution}{exo:b3:genomics:6}
$P(X < 8) = P(X \le 7) \approx P\bigl(Z < (7.5 - 15)/2.74\bigr) =
P(Z < -2.74) \approx 0.003$. À $30\times$, les deux allèles d'un
hétérozygote sont presque toujours vus de nombreuses fois, la couverture
est inégale (les régions riches en GC reçoivent moins de \omterm{def:b3:genomics:ngs}{lectures}, de
sorte que certains sites n'en voient que la moitié de la moyenne), et il
reste assez de \omterm{def:b3:genomics:ngs}{lectures} pour distinguer un vrai allèle des erreurs à
$10^{-3}$.
\end{solution}

\begin{solution}{exo:b3:genomics:7}
Toute \omterm{def:b3:genomics:ngs}{lecture} de \qty{150}{bp} prise à l'intérieur de la répétition est
identique quelle que soit la copie d'où elle vient, de sorte que
l'assembleur construit un seul sommet de \qty{6}{kb} avec $500$ chemins
entrants et $500$ sortants~; il ne peut pas dire quelle entrée va avec
quelle sortie, et l'\omterm{def:b3:genomics:assembly}{assemblage} se brise en $500$ lacunes, la répétition
n'étant présente qu'une fois avec une couverture $500$ fois la moyenne.
Des \omterm{def:b3:genomics:ngs}{lectures} plus longues que la répétition et ses flancs uniques des
deux côtés --- \qty{8}{kb} ou davantage --- franchissent chaque copie et
la résolvent.
\end{solution}

\begin{solution}{exo:b3:genomics:8}
$4.8\times 10^{7}/1000 = \num{48000}$ différences codantes~; un tiers,
environ \num{16000}, modifient la protéine.
\end{solution}

\begin{solution}{exo:b3:genomics:9}
Faux positifs attendus $10^{6}\times 5\times 10^{-8} = 0.05$. Un risque
relatif de $1.15$ sur une maladie à \qty{2}{\%} change quelques cas pour
mille~; mille personnes comptent une vingtaine de cas, bien trop peu
pour le voir (la puissance croît avec le nombre de cas et le carré de
l'effet). Pour un individu, $0.3$ point de pourcentage ne veut rien
dire. Mais l'allèle signale un gène dont un changement modeste
d'activité modifie le risque de maladie --- le gène, et sa voie, peuvent
être une cible médicamenteuse dont l'inhibition complète a un grand
effet.
\end{solution}

\begin{solution}{exo:b3:genomics:10}
Génétique des populations~: chez les bactéries, aux populations
énormes, la sélection élimine même les insertions un peu coûteuses~;
chez les mammifères, aux populations efficaces petites, la dérive laisse
s'accumuler un ADN non codant légèrement délétère --- étayé par la
corrélation inverse entre \omterm{prop:b3:genomics:sizes}{taille du génome} et taille de population selon
les lignées, miné par des exceptions. Structurale~: les génomes
eucaryotes sont envahis par des transposons que les bactéries, avec leur
biais vers la délétion et sans refuge méiotique pour les éléments
égoïstes, expulsent --- étayé par la corrélation entre \omterm{prop:b3:genomics:sizes}{taille du génome}
et contenu en transposons. Régulatrice~: un développement complexe
demande davantage d'ADN régulateur --- étayé par les \omterm{def:b3:genomics:comparative}{éléments non
codants conservés} autour des gènes du développement, miné par le fait
que \qty{8}{\%} seulement du génome humain montre la moindre sélection,
de sorte que la plus grande part de l'ADN non codant n'est pas
régulatrice.
\end{solution}

\begin{solution}{exo:b3:genomics:11}
Les départs des deux sortes tombent avec la densité totale $\rho = (N_{1} +
N_{2})/G$. Une \omterm{def:b3:genomics:ngs}{lecture} de longueur $L_{i}$ est la plus à droite de son
contig si aucune \omterm{def:b3:genomics:ngs}{lecture}, de l'une ou l'autre sorte, ne commence dans
les $L_{i} - T$ bases qui la suivent~: probabilité
$e^{-\rho(L_{i} - T)}$. D'où contigs $= N_{1}
e^{-\rho(L_{1} - T)} + N_{2} e^{-\rho(L_{2} - T)}$. Une longue \omterm{def:b3:genomics:ngs}{lecture}
est la plus à droite avec une probabilité exponentiellement plus faible
qu'une courte, et chaque longue \omterm{def:b3:genomics:ngs}{lecture} supprime en outre le risque
d'une lacune sur toute sa longueur~; le même nombre de bases en \omterm{def:b3:genomics:ngs}{lectures}
courtes ajoute beaucoup de points de départ mais chaque fenêtre protégée
est courte, de sorte que les longues \omterm{def:b3:genomics:ngs}{lectures} l'emportent.
\end{solution}

\begin{solution}{exo:b3:genomics:12}
Deux \omterm{def:b3:genomics:comparative}{duplications du génome entier} prédisent que le motif quadruple est
à l'échelle du génome~: des quatuors de segments chromosomiques
\omterm{def:b3:genomics:comparative}{paralogues} (paralogons) portant les mêmes familles de gènes dans le même
ordre, les duplications étant datées de la même époque pour toutes les
familles~; un complexe unique chez l'amphioxus et chez \emph{Ciona}, qui
ont divergé avant les événements~; et chez les lamproies, qui ont
divergé aux alentours, un état intermédiaire ou dérivé de façon
indépendante. Des duplications indépendantes du seul complexe Hox
prédisent des paralogons pour Hox seulement, des voisins d'histoires
différentes, et des dates de duplication qui diffèrent d'une famille à
l'autre. Les paralogons à l'échelle du génome et la datation commune,
tels qu'on les observe, plaident pour la \omterm{def:b3:genomics:comparative}{duplication du génome entier}.
\end{solution}

\begin{solution}{pb:b3:genomics:1}
\textbf{1.} $0.002\times 6\times 10^{8} = 1.2\times 10^{6}$ \omterm{def:b3:genomics:ngs}{lectures}~;
$c = 1.2\times 10^{6}\times 150/4.6\times 10^{6} \approx 39$.
\textbf{2.} $e^{-39} \approx 10^{-17}$~: en pratique aucune base non
couverte.
\textbf{3.} $\theta = 0.2$~; contigs $= 1.2\times 10^{6} e^{-31} \approx
0$~: un seul contig en théorie, des lacunes aux seules répétitions.
\textbf{4.} $c = \num{30000}\times 150/4.6\times 10^{6} = 0.98$~; non
couvert $e^{-0.98} = 0.38$, soit environ \qty{1.7}{Mb}~; contigs $= \num{30000}
\times e^{-0.78} \approx \num{13700}$.
\textbf{5.} Les sept opérons donnent des \omterm{def:b3:genomics:ngs}{lectures} identiques et
s'effondrent en un seul contig de \qty{5}{kb}, où entrent sept flancs
gauches uniques et d'où sortent sept flancs droits dont l'appariement
est inconnu~: $14$ extrémités de contigs, sept lacunes.
\textbf{6.} Environ \num{4600} gènes~; codant $0.88\times 4.6\times 10^{6}
= 4.0\times 10^{6}$ bp, gène moyen \qty{880}{bp}.
\textbf{7.} $0.998\times 6\times 10^{8}\times 150/3.2\times 10^{9}
\approx 28\times$.
\textbf{8.} Environ $14$ \omterm{def:b3:genomics:ngs}{lectures} par allèle.
\textbf{9.} $28\times 10^{-3}/3 \approx 0.009$ \omterm{def:b3:genomics:ngs}{lecture} montrant une base
fausse donnée --- la chance d'en avoir trois ou plus est d'environ
$10^{-7}$ --- et \qty{20}{\%} de $28$ font $5.6$ \omterm{def:b3:genomics:ngs}{lectures}~; une erreur
ne passe ni l'un ni l'autre test, tandis qu'un vrai allèle hétérozygote
est attendu dans $14$.
\textbf{10.} $3.2\times 10^{9}/1500 \approx 2.1\times 10^{6}$ sites
hétérozygotes~; en séquence codante $4.8\times 10^{7}/1500 \approx
\num{32000}$.
\textbf{11.} Une \omterm{def:b3:genomics:ngs}{lecture} de \qty{150}{bp} issue d'un Alu correspond
presque aussi bien à beaucoup des $1.1$ million de copies, de sorte
qu'elle est placée sur la mauvaise copie ou reçoit une faible confiance
d'alignement. Les différences entre copies ressemblent alors à des
variants hétérozygotes, et les vrais variants s'y perdent~: les \omterm{met:b3:genomics:pipeline}{appels
de variants} à l'intérieur d'un Alu sont d'ordinaire écartés, et ces
éléments sont les angles morts du séquençage à \omterm{def:b3:genomics:ngs}{lectures} courtes.
\textbf{12.} $1.2\times 10^{-8}\times 6.4\times 10^{9} \approx 77$
mutations nouvelles. Environ $14$ \omterm{def:b3:genomics:ngs}{lectures} devraient montrer le variant
chez l'enfant~; mais si le site d'un parent n'est couvert que par cinq
\omterm{def:b3:genomics:ngs}{lectures}, un allèle hétérozygote est manqué avec la probabilité
$2^{-5} = 3\,\%$ et un variant hérité est appelé à tort comme nouveau
--- les parents doivent donc être séquencés aussi profondément que
l'enfant.
\textbf{13.} $4.8\times 10^{7}\times 100/150 = 3.2\times 10^{7}$
\omterm{def:b3:genomics:ngs}{lectures}, vingt fois moins que le génome~: le coût de séquençage tombe
d'un facteur vingt, plus que ne coûte l'étape de capture.
\textbf{14.} $1.1\times 10^{6}\times 300 = 3.3\times 10^{8}$ bp, environ
\qty{10}{\%} du génome~; \qty{10}{\%} des \omterm{def:b3:genomics:ngs}{lectures}, soit quelque $6\times
10^{7}$.
\textbf{15.} $\num{60000}\times 1500 = 9\times 10^{7}$ bp~; $9\times
10^{7}/1.6\times 10^{10} = 0.56\,\%$.
\textbf{16.} $0.012\times 2.9\times 10^{9} = 3.5\times 10^{7}$
substitutions, $1.7\times 10^{7}$ par lignée~; taux $1.7\times
10^{7}/(2.9\times 10^{9}\times 6.5\times 10^{6}) = 9\times 10^{-10}$
par base et par an, $2.3\times 10^{-8}$ par génération de 25 ans ---
environ le double du taux généalogique de $1.2\times 10^{-8}$, écart
connu qui suggère des générations plus longues dans le passé ou une
séparation plus ancienne.
\textbf{17.} Quatre copies de \num{16000} feraient \num{64000}~;
\num{20000} subsistent, donc des \num{48000} copies supplémentaires
\num{4000} seulement survivent~: \qty{92}{\%} des doublons ont été
perdus.
\textbf{18.} Le ver et l'homme ont le même nombre de gènes~; la
complexité tient à l'usage qui en est fait --- l'épissage alternatif
multiplie un gène en milliers de protéines (\num{38016} pour
\emph{Dscam}), la régulation combine facteurs de transcription et
amplificateurs de façon propre à chaque type cellulaire, les \omterm{def:b3:rna-regulation:ncrna}{ARN non
codants} ajoutent des couches, et les protéines interagissent en réseaux.
\textbf{19.} Des éléments régulateurs (amplificateurs, promoteurs,
isolateurs), des gènes d'\omterm{def:b3:rna-regulation:ncrna}{ARN non codants}, des signaux d'épissage et de
localisation, des origines et des séquences structurales. On teste un
amplificateur candidat en le plaçant devant un gène rapporteur dans un
embryon transgénique et en regardant où le rapporteur est exprimé, puis
en supprimant l'élément du génome par CRISPR et en mesurant les gènes
voisins.
\textbf{20.} $0.05/10^{6} = 5\times 10^{-8}$~; à $p < 10^{-5}$, $10^{6}
\times 10^{-5} = 10$ faux positifs attendus.
\textbf{21.} Cote du non-porteur $0.009/0.991 = 0.00908$~; une copie,
cote $0.00954$, risque $0.945\,\%$~; deux copies, cote $0.01001$, risque
$0.99\,\%$.
\textbf{22.} Le score somme, sur les $200$ allèles, le nombre de copies
à risque que porte une personne, pondéré par le logarithme du rapport
des cotes de chaque allèle. Ce compte a pour moyenne
$200\times 2\times 0.3 = 120$ et pour écart type
$\sqrt{200\times 2\times 0.3\times 0.7} \approx 9$~; le centile
supérieur porte environ $21$ allèles de plus que la moyenne, et
$1.05^{21} \approx
2.8$~: près du triple de la cote d'une personne moyenne, à partir
d'allèles qui modifient chacun le risque de \qty{5}{\%}.
\textbf{23.} Les allèles d'un bloc de quelques dizaines de kilobases
sont hérités ensemble (déséquilibre de liaison), de sorte que n'importe
lequel montre la même association~; le SNP génotypé n'est que celui qui
figure sur la puce. Pour trouver le variant causal~: séquencer la région
chez les cas et les témoins, réduire l'ensemble par la statistique, puis
tester chaque candidat --- essais rapporteurs pour l'activité
amplificatrice de chaque allèle, édition de chaque allèle dans des
cellules et mesure de l'expression des gènes voisins, colocalisation
avec des signaux de caractères quantitatifs d'expression.
\textbf{24.} $1 - e^{-0.5} = 39\,\%$ du génome est lu. Cela donne un
échantillon direct d'une population à une date connue du passé --- les
fréquences alléliques avant les migrations et les métissages ultérieurs,
la longueur des segments néandertaliens (longs, parce que peu de
générations de recombinaison les avaient morcelés) et donc une date pour
le métissage --- ce que les génomes modernes ne peuvent qu'inférer par
des modèles.
\textbf{25.} Environ \num{13700} contigs dans l'essai~; $28\times$ par
génome haploïde, soit environ $14$ \omterm{def:b3:genomics:ngs}{lectures} par allèle~; quelque $77$
mutations nouvelles chez un enfant.
\end{solution}
